(1) 由 \(K=J/3\) 和 \(dK/d\tau=H\) 得 \(dJ/d\tau=3H\)。由于稳态净通量 \(F\) 常数，\(H=F/(4\pi)\) 不变。积分得 \(J(\tau)=3H\tau+C\)，顶部Eddington边界 \(J(0)=2H\) 给 \(C=2H\)，故 \(\boxed{J=3H(\tau+2/3)=3F(\tau+2/3)/(4\pi)}\)。这个 \(2/3\) 由指定的表面角分布闭合给出，并非所有真实大气都精确相同。

(2) 局域热平衡中积分Planck源函数为 \(B=\sigma_{\rm SB}T^4/\pi\)，令 \(J=B\) 并代入上式，得到 \(\boxed{T^4(\tau)=\frac34T_{\rm eff}^4(\tau+2/3)}\)。表面 \(\tau=0\) 给 \(T(0)=2^{-1/4}T_{\rm eff}\)，而 \(\tau=2/3\) 时 \(T=T_{\rm eff}\)。光球常取该光学深度，是灰色近似下的有效辐射层；不能把真实谱线形成深度一概固定为 \(2/3\)。

(3) 静力平衡和光学深度关系给 \(dP/d\tau=g/\kappa\)，按题设顶部零压积分为 \(\boxed{P=g\tau/\kappa}\)。由 \(T^4\propto\tau+2/3\) 得 \(d\ln T/d\tau=1/[4(\tau+2/3)]\)，而 \(d\ln P/d\tau=1/\tau\)，所以 \(\boxed{\nabla_{\rm rad}=\tau/[4(\tau+2/3)]<1/4}\)。单原子理想气体 \(\nabla_{\rm ad}=2/5\)，因 \(\nabla_{\rm rad}<\nabla_{\rm ad}\)，在这些常数不透明度、无其他能量源假设下满足Schwarzschild局部对流稳定判据。

(4) 直接微分温度律有 \(dT^4/d\tau=(3/4)T_{\rm eff}^4\)，代入扩散通量 \(F=(4\sigma_{\rm SB}/3)dT^4/d\tau\) 恰得 \(F=\sigma_{\rm SB}T_{\rm eff}^4\)，与最初给定通量一致。深层 \(\tau\gg1\) 时 \(T\propto\tau^{1/4}\)，扩散近似较好；近表面 \(\tau\lesssim1\) 辐射场明显各向异性，Eddington闭合仅是矩方法近似。若 \(\kappa\) 随温度变化、存在散射非LTE或外来照射，温度律和对流结论都需重算。
